\documentclass[12pt,a4paper]{article}
\usepackage[utf8]{inputenc}
\usepackage[T1]{fontenc}
\usepackage[spanish]{babel}
\usepackage{geometry}
\usepackage{xcolor}
\usepackage{listings}

\geometry{a4paper, margin=1in}

\lstset{
    backgroundcolor=\color{gray!5},
    basicstyle=\ttfamily\small,
    breaklines=true,
    keywordstyle=\color{blue},
    frame=single
}

\title{\textbf{Instrucciones para Consultar el Sistema RAG} \\ \large (Retrieval-Augmented Generation)}
\author{Prof. César Rodríguez \\ \small Curso: Seguridad Informática}
\date{\today}

\begin{document}

\maketitle

\section*{¿Qué es el Sistema RAG?}
El sistema RAG (\textit{Retrieval-Augmented Generation}) actúa como el ``segundo cerebro'' o la base de conocimiento de nuestra Suite de Auditoría. A diferencia de un buscador tradicional, entiende el \textbf{contexto semántico} de sus problemas. Si se topan con un error desarrollando sus módulos (ya sea de Git, Python, Docker o del orquestador), es muy probable que el sistema ya conozca la solución basada en los errores previos del curso.

\section*{¿Cómo realizar una consulta?}
Para consultar la base de datos (impulsada por ChromaDB), utilizaremos el script \texttt{knowledge\_db.py} desde la terminal. 

La estructura básica del comando es:
\begin{lstlisting}[language=bash]
python knowledge_db.py -c <coleccion> -q 'Tu pregunta en lenguaje natural'
\end{lstlisting}

\textbf{Parámetros:}
\begin{itemize}
    \item \texttt{-c <coleccion>}: Especifica la colección temática a consultar (ej. \texttt{auditoria}, \texttt{python}, \texttt{git}).
    \item \texttt{-q '<texto>'}: Es la bandera de consulta (\textit{query}) seguida de su pregunta o descripción del problema entre comillas.
\end{itemize}

\section*{Ejemplos de Consultas}
No necesitan recordar palabras exactas o códigos de error. Escriban con naturalidad describiendo su problema:

\textbf{Ejemplo 1: Problemas de ejecución (Rate Limits, bloqueos, etc.)}
\begin{lstlisting}[language=bash]
python knowledge_db.py -c auditoria -q 'por que mi script de google dorks se detiene o no saca resultados'
\end{lstlisting}

\textbf{Ejemplo 2: Errores devueltos por la suite principal}
\begin{lstlisting}[language=bash]
python knowledge_db.py -c auditoria -q 'el orquestador me da un error de validacion JSON en la llave status'
\end{lstlisting}

\vspace{1cm}
\noindent \textit{Nota: Si el RAG no tiene la respuesta, recuerden documentar la solución una vez que la encuentren para alimentar la base de datos, tal como se especifica en la \textbf{Guía del Sistema RAG} (\texttt{guia\_rag.pdf}).}

\end{document}